\documentclass[11pt,a4paper]{article}

% ─── Packages ────────────────────────────────────────────────────────────────
\usepackage[utf8]{inputenc}
\usepackage[T1]{fontenc}
\usepackage{lmodern}
\usepackage[margin=1in]{geometry}
\usepackage{graphicx}
\usepackage{booktabs}
\usepackage{longtable}
\usepackage{enumitem}
\usepackage{hyperref}
\usepackage{xcolor}
\usepackage{listings}
\usepackage{amsmath}
\usepackage{caption}
\usepackage{float}
\usepackage{fancyhdr}
\usepackage{titlesec}

% ─── Code listing style ─────────────────────────────────────────────────────
\definecolor{codebg}{HTML}{F5F5F5}
\definecolor{codeframe}{HTML}{CCCCCC}
\definecolor{keyword}{HTML}{0000FF}
\definecolor{comment}{HTML}{008000}
\definecolor{string}{HTML}{A31515}

\lstdefinestyle{pythonstyle}{
    language=Python,
    basicstyle=\ttfamily\footnotesize,
    backgroundcolor=\color{codebg},
    frame=single,
    rulecolor=\color{codeframe},
    keywordstyle=\color{keyword}\bfseries,
    commentstyle=\color{comment}\itshape,
    stringstyle=\color{string},
    showstringspaces=false,
    breaklines=true,
    breakatwhitespace=true,
    tabsize=4,
    numbers=left,
    numberstyle=\tiny\color{gray},
    numbersep=8pt,
    xleftmargin=2em,
    framexleftmargin=1.5em,
    captionpos=b,
}

\lstdefinestyle{tomlstyle}{
    basicstyle=\ttfamily\footnotesize,
    backgroundcolor=\color{codebg},
    frame=single,
    rulecolor=\color{codeframe},
    keywordstyle=\color{keyword}\bfseries,
    commentstyle=\color{comment}\itshape,
    stringstyle=\color{string},
    showstringspaces=false,
    breaklines=true,
    tabsize=4,
    numbers=left,
    numberstyle=\tiny\color{gray},
    numbersep=8pt,
    xleftmargin=2em,
    framexleftmargin=1.5em,
    captionpos=b,
}

\lstdefinestyle{shellstyle}{
    basicstyle=\ttfamily\footnotesize,
    backgroundcolor=\color{codebg},
    frame=single,
    rulecolor=\color{codeframe},
    showstringspaces=false,
    breaklines=true,
    tabsize=4,
    xleftmargin=1em,
    framexleftmargin=0.5em,
    captionpos=b,
}

\lstset{style=pythonstyle}

% ─── Header/Footer ──────────────────────────────────────────────────────────
\pagestyle{fancy}
\fancyhf{}
\fancyhead[L]{\small\textsc{BCI Interface --- AFE Headstage}}
\fancyhead[R]{\small\textsc{CI Guardrails Report}}
\fancyfoot[C]{\thepage}
\renewcommand{\headrulewidth}{0.4pt}

% ─── Title ───────────────────────────────────────────────────────────────────
\title{%
    \textbf{Hermetic CI Guardrails for Neural AFE PCB\\Design Rule Check Pipeline}\\[0.5em]
    \large Design, Implementation, and Verification Report\\[0.3em]
    \normalsize BCI Interface Project --- AFE Headstage Real
}
\author{%
    Hardware Verification Team\\
    \texttt{hardware/afe-headstage-real/}
}
\date{February 26, 2026}

\begin{document}

\maketitle

\begin{abstract}
This report documents the design, iterative implementation, and verification of a
hermetic continuous integration (CI) governance framework for the neural analog
front-end (AFE) headstage PCB's Design Rule Check (DRC) pipeline. The system
enforces structured parsing of KiCad~9 DRC reports with regression thresholds,
canonical source file formatting, single-source-of-truth regex invariants,
a tiered static analysis policy via Ruff, and a standalone CI policy guard script
that prevents ``green by suppression'' --- the common failure mode where lint
configurations silently hide real defects. The final system comprises 1{,}425~lines
across four artifacts, passes a five-stage CI verification chain, and enforces
policy through 27~automated tests covering parser correctness, mutation detection,
format canonicalization, and configuration governance.
\end{abstract}

\tableofcontents
\newpage

% ══════════════════════════════════════════════════════════════════════════════
\section{Introduction}
\label{sec:intro}

\subsection{Motivation}

Brain--computer interface (BCI) headstage PCBs demand high-reliability electrical
verification. KiCad~9's built-in Design Rule Check (DRC) produces text reports
enumerating violations across clearance, unconnected nets, via connectivity,
silkscreen overlaps, and library footprint mismatches. Manually reviewing these
reports is error-prone and does not scale across design revisions.

A common failure mode in hardware CI is \emph{green by suppression}: lint or DRC
configurations are progressively weakened (via broad ignore rules, file exclusions,
or suppressed violation families) until the pipeline reports zero errors --- not
because defects are fixed, but because they are hidden. This is particularly
dangerous for safety-critical categories like unconnected pads on signal nets.

\subsection{Objectives}

This work establishes a CI governance framework with four properties:

\begin{enumerate}[label=\textbf{P\arabic*}]
    \item \textbf{Structured DRC Parsing:} Parse KiCad~9 \texttt{.rpt} files into
          typed histograms with cross-checked section totals and regression thresholds.
    \item \textbf{Invariant Enforcement:} Guarantee canonical source file headers,
          single-source-of-truth regex definitions, and AST-verified module structure.
    \item \textbf{Tiered Static Analysis:} Apply Ruff lint/format rules with a
          tiered policy --- strict for CI-critical files, cosmetic-only tolerances
          for legacy generators --- with no suppression of correctness rules (F-family).
    \item \textbf{Configuration Governance:} A standalone policy guard that
          structurally parses \texttt{pyproject.toml} via TOML and rejects any
          configuration that could weaken F-rule enforcement, including bypass
          vectors such as \texttt{exclude}, \texttt{extend\-exclude},
          \texttt{extend\-per\-file\-ignores}, and \texttt{select} manipulation.
\end{enumerate}

\subsection{Toolchain}

\begin{table}[H]
\centering
\caption{Development and CI toolchain versions.}
\label{tab:toolchain}
\begin{tabular}{@{}lll@{}}
\toprule
\textbf{Tool} & \textbf{Version} & \textbf{Role} \\
\midrule
Python       & 3.14.3  & Runtime, \texttt{tomllib} (stdlib), \texttt{ast} module \\
Ruff         & 0.15.2  & Linting (E/F/I families), import sorting, formatting \\
pytest       & 9.0.2   & Test framework (27 tests) \\
KiCad        & 9.x     & DRC report generator (upstream) \\
\texttt{tomllib} & stdlib & TOML parsing for policy guards \\
\bottomrule
\end{tabular}
\end{table}

% ══════════════════════════════════════════════════════════════════════════════
\section{Architecture}
\label{sec:arch}

The system consists of four artifacts totaling 1{,}425~lines of code:

\begin{table}[H]
\centering
\caption{Artifact inventory.}
\label{tab:artifacts}
\begin{tabular}{@{}llrl@{}}
\toprule
\textbf{File} & \textbf{Role} & \textbf{Lines} & \textbf{Category} \\
\midrule
\texttt{parse\_drc.py}       & DRC report parser       & 494 & Parser \\
\texttt{test\_parse\_drc.py} & Test suite (27 tests)   & 750 & Verification \\
\texttt{ci\_lint\_policy.py} & CI policy guard         & 158 & Governance \\
\texttt{pyproject.toml}      & Ruff configuration      &  23 & Configuration \\
\midrule
\multicolumn{2}{@{}l}{\textbf{Total}} & \textbf{1{,}425} & \\
\bottomrule
\end{tabular}
\end{table}

The CI execution order is strict and non-negotiable:

\begin{lstlisting}[style=shellstyle,caption={CI execution chain (5 stages).},label={lst:ci-chain}]
python ci_lint_policy.py          # Stage 1: Policy guard
python -m compileall . -q         # Stage 2: Syntax check
ruff check .                      # Stage 3: Lint
ruff format --check *.py          # Stage 4: Format check
pytest -q test_parse_drc.py       # Stage 5: Tests (27)
\end{lstlisting}

Stage~1 \emph{must} precede Stage~3: the policy guard validates that the Ruff
configuration itself has not been weakened before Ruff is allowed to run.

% ══════════════════════════════════════════════════════════════════════════════
\section{DRC Report Parser}
\label{sec:parser}

\subsection{KiCad 9 Report Structure}

KiCad~9 DRC reports contain three independently counted sections:

\begin{enumerate}
    \item \texttt{** Found N DRC violations **} --- clearance, silkscreen,
          via dangling, library footprint issues
    \item \texttt{** Found M unconnected pads **} --- missing net connections
    \item \texttt{** Found P Footprint errors **} --- component footprint mismatches
\end{enumerate}

Each violation is a multi-line block beginning with a bracketed type key
(e.g., \texttt{[clearance]}, \texttt{[unconnected\_items]}) followed by
indented body lines containing coordinate/pad/net annotations.

\subsection{Parser Design}

\texttt{parse\_drc.py} (494~lines) implements:

\begin{itemize}
    \item \textbf{Block-scoped parsing:} Violations are only counted when their
          \texttt{[type\_key]} header is encountered. Pad lines outside a
          recognized block are ignored (mutation-tested, see \S\ref{sec:mutation}).
    \item \textbf{Cross-checked totals:} Parsed histogram counts are reconciled
          against KiCad's declared section totals.
    \item \textbf{Net extraction:} Net names are extracted from bracket annotations
          (\texttt{[NET\_NAME]}) with layer-span filtering (\texttt{F.Cu - B.Cu}
          is not a net name).
    \item \textbf{Unconnected identity extraction:} For \texttt{[unconnected\_items]}
          blocks, structured \texttt{(ref, pad, net, layer)} tuples are extracted
          via \texttt{PAD\_LINE\_RE}.
\end{itemize}

\subsection{Key Regex: \texttt{PAD\_LINE\_RE}}

A single module-level compiled regex handles pad identity extraction:

\begin{lstlisting}[caption={\texttt{PAD\_LINE\_RE} --- single source of truth (line 40).},label={lst:pad-re}]
PAD_LINE_RE = re.compile(
    r"""
    @\(\s*
        (?P<x>-?\d+(?:\.\d+)?)\s*mm\s*,\s*
        (?P<y>-?\d+(?:\.\d+)?)\s*mm
    \s*\)\s*:\s*
    Pad\s+(?P<pad>[A-Za-z0-9_.-]+)\s*
    \[\s*(?P<net>[^\]]+?)\s*\]\s*
    of\s+(?P<ref>[A-Za-z0-9_.-]+)\s+
    on\s+(?P<layer>[A-Za-z0-9. ]+?)\s*,?\s*$
    """,
    re.VERBOSE,
)
\end{lstlisting}

The tolerant layer pattern \texttt{[A-Za-z0-9.\textbackslash{} ]+?} and optional
trailing comma handle KiCad format drift across minor versions.

\subsection{Regression Thresholds}

The parser enforces version-tagged ceiling constants:

\begin{itemize}
    \item \texttt{UNCONNECTED\_CEILING} --- maximum allowed unconnected pad count
    \item \texttt{LIB\_FOOTPRINT\_ISSUES\_CEILING} --- maximum library footprint issues
    \item \texttt{VIA\_DANGLING\_ALLOWLIST} --- frozen set of \texttt{(x, y, net)} tuples
          (empty since v35: all stitch vias connect via GND pours)
    \item \texttt{EXPECTED\_DEFAULT\_CLEARANCE} --- board clearance sanity bound
\end{itemize}

% ══════════════════════════════════════════════════════════════════════════════
\section{Test Suite}
\label{sec:tests}

The test suite (\texttt{test\_parse\_drc.py}, 750~lines) contains 27~tests
organized into six categories:

{\small
\begin{longtable}{@{}p{5.2cm}p{1.6cm}p{5.8cm}@{}}
\caption{Test suite inventory (27 tests, all passing).}
\label{tab:tests} \\
\toprule
\textbf{Test Name} & \textbf{Cat.} & \textbf{What It Verifies} \\
\midrule
\endfirsthead
\toprule
\textbf{Test Name} & \textbf{Cat.} & \textbf{What It Verifies} \\
\midrule
\endhead
\texttt{test\_basic\_counts} & Parser & Section total reconciliation \\
\texttt{test\_net\_extraction} & Parser & Net name extraction, layer filter \\
\texttt{test\_unconnected\_allowlist} & Parser & Power nets pass allowlist \\
\texttt{test\_signal\_net\_unconn\_flagged} & Parser & Signal nets flagged \\
\texttt{test\_count\_mismatch\_detected} & Parser & Declared vs.\ parsed mismatch \\
\texttt{test\_empty\_report\_raises} & Parser & Empty file $\to$ RuntimeError \\
\texttt{test\_whitespace\_only\_raises} & Parser & Whitespace-only $\to$ RuntimeError \\
\texttt{test\_missing\_drc\_header\_raises} & Parser & No header $\to$ RuntimeError \\
\texttt{test\_zero\_violations} & Parser & Clean report handling \\
\texttt{test\_multi\_nets\_in\_violation} & Parser & Multi-net block \\
\texttt{test\_baseline\_consts\_positive} & Sanity & Ceiling $\geq 0$, version tag \\
\texttt{test\_via\_dangle\_allow\_frozen} & Sanity & Type/shape invariant \\
\texttt{test\_via\_dangle\_allow\_empty} & Sanity & v35 invariant \\
\texttt{test\_via\_dangle\_non\_allow\_det} & Parser & Non-allowlisted flagged \\
\texttt{test\_via\_dangle\_wrong\_net\_rej} & Parser & Coordinate match, net mismatch \\
\texttt{test\_lib\_fp\_counted\_separately} & Parser & Separate histogram bucket \\
\texttt{test\_default\_clearance\_sane} & Sanity & $0.05 \leq c \leq 0.5$ \\
\texttt{test\_unconn\_extracts\_expected} & Fixture & Real fixture identity check \\
\texttt{test\_unconn\_format\_drift\_\dots} & Robust. & 6 format variants \\
\texttt{test\_strict\_unconn\_empty\_\dots} & Fixture & Empty fixture $\to \emptyset$ \\
\texttt{test\_strict\_unconn\_rejects\_mut} & Mutation & Injected pad detected \\
\texttt{test\_pad\_outside\_block\_not\_ct} & Mutation & Block scope enforced \\
\texttt{test\_future\_import\_unique} & Format & \texttt{\_\_future\_\_} line 1 \\
\texttt{test\_single\_pad\_line\_re\_def} & Format & One \texttt{PAD\_LINE\_RE =} \\
\texttt{test\_docstring\_after\_future} & Format & AST: stmt[1] is docstring \\
\texttt{test\_no\_f\_suppress\_legacy} & Gov. & No F* in per-file-ignores \\
\texttt{test\_exclude\_knobs\_forbidden} & Gov. & No exclude/extend-exclude keys \\
\bottomrule
\end{longtable}
}

\subsection{Mutation Testing}
\label{sec:mutation}

Two tests implement targeted mutation testing without a mutation framework:

\begin{enumerate}
    \item \textbf{\texttt{test\_pad\_line\_outside\_block\_not\_counted}}: Injects a
          valid \texttt{Pad} line \emph{outside} any \texttt{[unconnected\_items]}
          block into a real fixture file. Asserts the parser does not count it.
          Then injects the same line \emph{inside} a block and asserts it is counted.
    \item \textbf{\texttt{test\_strict\_unconnected\_rejects\_mutation\_on\_empty\_fixture}}:
          Takes a zero-unconnected fixture, injects a fake unconnected pad entry,
          and asserts the resulting identity set is non-empty (i.e., the mutation
          is detected).
\end{enumerate}

\subsection{Format Drift Robustness}

\texttt{test\_unconnected\_pad\_line\_format\_drift\_variants\_parse\_same} tests
six spacing/punctuation variants of the same pad line:

\begin{lstlisting}[caption={Format drift test variants.},label={lst:drift}]
"@(123.5000 mm, 100.5000 mm): Pad M16 [GND] of U1 on F.Cu"
"@( 123.5000mm,100.5000mm ):  Pad   M16   [GND]  of   U1  on  F.Cu "
"@(123.5000 mm, 100.5000 mm): Pad M16 [GND] of U1 on F.Cu,"
"    @(123.5000 mm, 100.5000 mm): Pad M16 [GND] of U1 on F.Cu"
"@(123.5000 mm, 100.5000 mm): Pad M16-2 [GND] of U1-1 on F.Cu,"
"@(123.5000 mm, 100.5000 mm): Pad M16.2 [GND] of U1.1 on F.Cu,"
\end{lstlisting}

All six must parse to the same identity tuple, ensuring the regex survives
minor KiCad formatting changes.

% ══════════════════════════════════════════════════════════════════════════════
\section{CI Lint Policy}
\label{sec:lint-policy}

\subsection{Tiered Ruff Configuration}

The Ruff linter is configured with a tiered policy in \texttt{pyproject.toml}:

\begin{lstlisting}[style=tomlstyle,caption={Ruff configuration (pyproject.toml).},label={lst:pyproject}]
[tool.ruff]
target-version = "py311"
line-length = 100

[tool.ruff.lint]
select = ["E", "F", "I"]
ignore = ["E402", "E741"]

[tool.ruff.lint.per-file-ignores]
"fill_zones.py" = ["E501"]
"gen_*.py" = ["E501", "E401", "I001"]
"verify_*.py" = ["E501", "E401", "I001"]
"test_*.py" = ["E501"]

[tool.ruff.lint.isort]
known-local-folder = ["parse_drc"]
\end{lstlisting}

\textbf{Design principle:} Legacy generator/verifier scripts (\texttt{gen\_*.py},
\texttt{verify\_*.py}) are allowed only \emph{cosmetic} tolerances (long lines,
import ordering). No F-family rules (unused imports, unused variables, f-string
placeholders) are suppressed anywhere. This forces actual cleanup of dead code
rather than hiding it.

\subsection{Policy Guard Script}
\label{sec:guard}

\texttt{ci\_lint\_policy.py} (158~lines) is the authoritative CI governance
mechanism. It parses \texttt{pyproject.toml} via \texttt{tomllib} and enforces
five guards:

\begin{table}[H]
\centering
\caption{Policy guard inventory.}
\label{tab:guards}
\begin{tabular}{@{}clp{7.5cm}@{}}
\toprule
\textbf{Guard} & \textbf{Name} & \textbf{What It Enforces} \\
\midrule
0a & Select union & $\texttt{select} \cup \texttt{extend-select}$ must include
     \texttt{"F"} or \texttt{"ALL"} \\
0b & Exclude ban & \texttt{exclude}, \texttt{extend-exclude},
     \texttt{force-exclude} forbidden under both \texttt{[tool.ruff]} and
     \texttt{[tool.ruff.lint]} \\
0c & Extend-per-file ban & \texttt{extend-per-file-ignores} forbidden \\
1  & Global F* ban & \texttt{ignore} and \texttt{extend-ignore} cannot contain
     any code starting with \texttt{"F"} \\
2  & Legacy per-file F* ban & Per-file-ignores keys whose basename starts with
     \texttt{gen\_} or \texttt{verify\_} cannot contain any \texttt{F*} code;
     invalid value types hard-fail \\
\bottomrule
\end{tabular}
\end{table}

\subsubsection{Basename Matching}

Per-file-ignores keys are TOML strings that may contain path prefixes or glob
patterns (e.g., \texttt{"**/verify\_*.py"}, \texttt{"scripts/gen\_pcb.py"}).
The guard extracts the basename before prefix matching:

\begin{lstlisting}[caption={Basename extraction for glob-safe legacy key matching.},label={lst:basename}]
def _is_legacy_key(key: str) -> bool:
    last = key.replace("\\", "/").split("/")[-1]
    return any(last.startswith(p) for p in LEGACY_KEY_PREFIXES)
\end{lstlisting}

This prevents bypass via path-qualified keys.

\subsubsection{Union Select Check}

Guard~0a evaluates the \emph{combined} effective selection, not each field
independently. This correctly accepts configurations like:

\begin{lstlisting}[style=tomlstyle]
select = ["E", "I"]
extend-select = ["F"]
\end{lstlisting}

\noindent while rejecting:

\begin{lstlisting}[style=tomlstyle]
select = ["E", "I"]
# No extend-select => F not in union => FAIL
\end{lstlisting}

\subsubsection{Type Safety}

The \texttt{\_iter\_codes()} helper handles three TOML value shapes:

\begin{itemize}
    \item \texttt{str} --- single code (returns \texttt{([val], True)})
    \item \texttt{list[str]} --- code list (returns \texttt{(list(val), True)})
    \item anything else --- returns \texttt{([], False)}, triggering a hard
          failure (exit~2) rather than silent skip
\end{itemize}

\subsection{Exit Code Semantics}

\begin{table}[H]
\centering
\caption{\texttt{ci\_lint\_policy.py} exit codes.}
\label{tab:exit}
\begin{tabular}{@{}cl@{}}
\toprule
\textbf{Code} & \textbf{Meaning} \\
\midrule
0 & Policy clean \\
1 & Policy violation (forbidden suppression detected) \\
2 & Structural error (missing file, invalid TOML shape) \\
\bottomrule
\end{tabular}
\end{table}

% ══════════════════════════════════════════════════════════════════════════════
\section{Threat Model}
\label{sec:threats}

Table~\ref{tab:threats} enumerates known bypass vectors and their mitigations:

\begin{longtable}{@{}p{4cm}p{2cm}p{7cm}@{}}
\caption{Bypass threat model and mitigations.}
\label{tab:threats} \\
\toprule
\textbf{Bypass Vector} & \textbf{Guard} & \textbf{Mitigation} \\
\midrule
\endfirsthead
\toprule
\textbf{Bypass Vector} & \textbf{Guard} & \textbf{Mitigation} \\
\midrule
\endhead
Add \texttt{F401} to global ignore &
    Guard~1 & Scans \texttt{ignore} and \texttt{extend-ignore} for F* prefix \\
\addlinespace
Add F841 to legacy per-file-ignores &
    Guard~2 & Scans legacy-bucket values for F* prefix \\
\addlinespace
Path-qualified key (e.g.\ \texttt{**/gen\_*.py}) &
    Guard~2 & Basename extraction via \texttt{split("/")} \\
\addlinespace
Remove \texttt{F} from \texttt{select} &
    Guard~0a & Union check on select $\cup$ extend-select \\
\addlinespace
Add \texttt{exclude} array &
    Guard~0b & All exclude variants banned outright \\
\addlinespace
Add \texttt{extend-exclude} &
    Guard~0b & Checked alongside exclude/force-exclude \\
\addlinespace
Add \texttt{extend-per-file-ignores} &
    Guard~0c & Key banned outright \\
\addlinespace
Duplicate \texttt{PAD\_LINE\_RE} &
    Test & Regex count test (single definition) \\
\addlinespace
Pad line outside block &
    Test & Block-scope mutation test \\
\addlinespace
Move docstring after code &
    Test & AST-based position test \\
\addlinespace
Non-str value in per-file-ignores &
    Guard~2 & \texttt{\_iter\_codes} returns \texttt{ok=False} $\to$ exit~2 \\
\addlinespace
Run ruff before policy guard &
    CI order & Policy is Stage~1; Ruff is Stage~3 \\
\bottomrule
\end{longtable}

\subsection{Residual Risks}

\begin{enumerate}
    \item \textbf{Inline \texttt{\# noqa: F401} comments:} Per-line suppressions
          within source files are not currently blocked by the policy guard.
          Mitigation: code review + future \texttt{grep}-based guard.
    \item \textbf{\texttt{[tool.ruff] src} manipulation:} Changing the source
          root could affect which files Ruff sees. Low risk given the flat
          directory structure.
    \item \textbf{CI script itself being skipped:} If the GitHub Actions workflow
          is modified to remove Stage~1, governance is lost. Mitigation: the pytest
          backstop tests (\texttt{test\_no\_f\_rule\_suppression...},
          \texttt{test\_exclude\_knobs\_forbidden}) provide a secondary check
          during Stage~5.
\end{enumerate}

% ══════════════════════════════════════════════════════════════════════════════
\section{Legacy Code Cleanup}
\label{sec:cleanup}

Prior to establishing the tiered policy, legacy files contained numerous
F-family violations that would have been suppressed under a naive configuration.
These were fixed, not suppressed:

\begin{table}[H]
\centering
\caption{Legacy file cleanup summary.}
\label{tab:cleanup}
\begin{tabular}{@{}llrl@{}}
\toprule
\textbf{File} & \textbf{Rule} & \textbf{Count} & \textbf{Fix} \\
\midrule
\texttt{gen\_pcb.py}          & F541 & 31 & Removed empty f-string prefixes \\
\texttt{gen\_pcb.py}          & ---  &  1 & Deleted dead \texttt{\_R\_PAD\_HW\_X} symbol \\
\texttt{gen\_schematic.py}    & F401 &  1 & Removed unused \texttt{textwrap} import \\
\texttt{gen\_schematic.py}    & F841 &  3 & Removed 3 dead variables \\
\texttt{gen\_schematic\_v3.py}& F401 &  1 & Removed unused \texttt{math} import \\
\texttt{verify\_schematic.py} & F841 &  1 & Removed dead \texttt{net\_code} variable \\
\texttt{verify\_pcb\_guards.py}& F541 & 5 & Removed empty f-string prefixes \\
\texttt{fill\_zones.py}       & F541 &  1 & Removed empty f-string prefix \\
\midrule
\multicolumn{2}{@{}l}{\textbf{Total fixes}} & \textbf{44} & \\
\bottomrule
\end{tabular}
\end{table}

% ══════════════════════════════════════════════════════════════════════════════
\section{Verification Results}
\label{sec:results}

\subsection{Final CI Chain Output}

The complete five-stage CI chain produces the following output on the final
codebase state:

\begin{lstlisting}[style=shellstyle,caption={Final CI verification output (Feb 26, 2026).},label={lst:results},basicstyle=\ttfamily\scriptsize]
$ python ci_lint_policy.py
POLICY: 0

$ python -m compileall . -q
COMPILE: 0

$ ruff check .
All checks passed!
RUFF: 0

$ ruff format --check *.py
3 files already formatted
FORMAT: 0

$ pytest test_parse_drc.py -v --tb=no
::test_pad_line_outside_block_not_counted    PASSED
::test_future_import_first_and_unique        PASSED
::test_single_pad_line_re_definition         PASSED
::test_docstring_immediately_after_future..  PASSED
::test_no_f_rule_suppression_in_legacy_...   PASSED
::test_exclude_knobs_forbidden               PASSED
::test_basic_counts                          PASSED
::test_net_extraction                        PASSED
::test_unconnected_allowlist                 PASSED
::test_signal_net_unconnected_flagged        PASSED
::test_count_mismatch_detected               PASSED
::test_empty_report_raises                   PASSED
::test_whitespace_only_raises                PASSED
::test_missing_drc_header_raises             PASSED
::test_zero_violations                       PASSED
::test_multiple_nets_in_one_violation        PASSED
::test_baseline_constants_are_positive       PASSED
::test_via_dangling_allowlist_is_frozen      PASSED
::test_via_dangling_allowlist_is_empty       PASSED
::test_via_dangling_non_allowlisted_detected PASSED
::test_via_dangling_wrong_net_rejected       PASSED
::test_lib_footprint_counted_separately      PASSED
::test_expected_default_clearance_is_sane    PASSED
::test_unconnected_extracts_expected_...     PASSED
::test_unconnected_pad_line_format_drift_..  PASSED
::test_strict_unconnected_matches_empty_..   PASSED
::test_strict_unconnected_rejects_mutation.. PASSED
=============== 27 passed in 0.04s ================
PYTEST: 0
\end{lstlisting}

\subsection{Summary Metrics}

\begin{table}[H]
\centering
\caption{Final verification summary.}
\label{tab:summary}
\begin{tabular}{@{}lr@{}}
\toprule
\textbf{Metric} & \textbf{Value} \\
\midrule
Total source lines             & 1{,}425 \\
Policy guards                  & 5 (Guards 0a, 0b, 0c, 1, 2) \\
Bypass vectors mitigated       & 12 \\
Tests passing                  & 27 / 27 \\
Ruff violations                & 0 \\
Format drift                   & 0 files \\
Legacy F-violations fixed      & 44 \\
Legacy F-violations suppressed & 0 \\
CI stages                      & 5 \\
CI execution time              & $< 1$~s \\
\bottomrule
\end{tabular}
\end{table}

% ══════════════════════════════════════════════════════════════════════════════
\section{Iterative Development Process}
\label{sec:process}

The system was developed over approximately 12 iterative rounds, each driven
by a specific failure mode or loophole discovery:

\begin{enumerate}
    \item \textbf{Parser hardening:} Strict unconnected identity extraction with
          block-scoped tracking. Mutation-tested with fixture injection.
    \item \textbf{Import/header canonicalization:} Canonical header format
          (\texttt{from \_\_future\_\_ import annotations} on line~1, docstring
          on line~2). Enforced via AST-based test.
    \item \textbf{Single \texttt{PAD\_LINE\_RE}:} Discovered and removed duplicate
          definition (module-level and function-level). Enforced via
          \texttt{re.findall} count test.
    \item \textbf{Ruff configuration + legacy cleanup:} Installed Ruff, created
          tiered policy. Fixed 44 F-family violations across 6 legacy files.
    \item \textbf{Policy tightening:} Rejected initial broad F-rule suppression.
          Restricted legacy per-file-ignores to cosmetic-only codes.
    \item \textbf{CI policy guard (v1):} String-based TOML checking. Upgraded to
          \texttt{tomllib}-based structural parsing.
    \item \textbf{Global ignore/extend-ignore guard:} Closed loophole where F*
          codes could be added to global ignore list.
    \item \textbf{Select union check:} Fixed false rejection of valid
          \texttt{select} + \texttt{extend-select} combinations.
    \item \textbf{Exclude ban:} Banned all exclude variants (\texttt{exclude},
          \texttt{extend-exclude}, \texttt{force-exclude}) under both
          \texttt{[tool.ruff]} and \texttt{[tool.ruff.lint]}.
    \item \textbf{Basename matching:} Fixed legacy key matching to use basename
          of glob pattern, preventing path-prefix bypass.
    \item \textbf{Extend-per-file-ignores ban:} Banned the key outright to
          prevent a secondary per-file-ignores channel.
    \item \textbf{Pytest backstop tests:} Added \texttt{test\_exclude\_knobs\_forbidden}
          as belt-and-suspenders behind the CI policy script.
\end{enumerate}

% ══════════════════════════════════════════════════════════════════════════════
\section{Conclusion}
\label{sec:conclusion}

This work demonstrates that ``green CI'' for hardware DRC pipelines requires
not just parsing and testing, but \emph{governance of the CI configuration itself}.
The primary contribution is a standalone policy guard (\texttt{ci\_lint\_policy.py})
that structurally validates the linter configuration before the linter runs,
closing 12 known bypass vectors including path-glob evasion, exclude knobs,
global suppression, and select manipulation.

The system achieves zero suppressed F-family violations across 44 previously
hidden defects in legacy code, enforced by 27 automated tests executing in
under one second. The five-stage CI chain is ordered such that configuration
governance precedes all downstream checks, establishing the invariant that
a passing pipeline reflects genuine code quality rather than configuration
erosion.

\subsection{Future Work}

\begin{itemize}
    \item Inline \texttt{\# noqa} scanning to detect per-line F-rule suppression
    \item GitHub Actions workflow integration with branch protection rules
    \item Coverage measurement and enforcement for \texttt{parse\_drc.py}
    \item Extension of policy guards to additional rule families (e.g., B-family
          from flake8-bugbear) as they are enabled
    \item Automated DRC report generation via KiCad CLI in CI
\end{itemize}

% ══════════════════════════════════════════════════════════════════════════════
\appendix

\section{Complete Policy Guard Source}
\label{app:policy}

\begin{lstlisting}[caption={\texttt{ci\_lint\_policy.py} (158 lines).},label={lst:policy-full}]
from __future__ import annotations

import sys
from pathlib import Path

try:
    import tomllib  # py>=3.11
except ModuleNotFoundError:  # pragma: no cover
    import tomli as tomllib  # type: ignore[no-redef]


LEGACY_KEY_PREFIXES = ("gen_", "verify_")
FORBIDDEN_PREFIX = "F"


def _is_legacy_key(key: str) -> bool:
    last = key.replace("\\", "/").split("/")[-1]
    return any(last.startswith(p) for p in LEGACY_KEY_PREFIXES)


def _iter_codes(val: object) -> tuple[list[str], bool]:
    """Return (codes, ok_type)."""
    if isinstance(val, str):
        return [val], True
    if isinstance(val, list):
        if all(isinstance(x, str) for x in val):
            return list(val), True
        return [x for x in val if isinstance(x, str)], False
    return [], False


def _strip_code(code: str) -> str:
    return code.strip().strip('"').strip("'")


def main() -> int:
    pyproject = Path(__file__).with_name("pyproject.toml")
    if not pyproject.exists():
        print(f"ERROR: {pyproject} not found", file=sys.stderr)
        return 2

    data = tomllib.loads(pyproject.read_text(encoding="utf-8"))

    tool = data.get("tool", {})
    ruff = tool.get("ruff", {})
    lint = ruff.get("lint", {})

    # Guard 0a: union of select + extend-select must include F (or ALL).
    effective_select: set[str] = set()
    for sel_field in ("select", "extend-select"):
        sel_val = lint.get(sel_field)
        if sel_val is None:
            continue
        codes, ok = _iter_codes(sel_val)
        if not ok:
            print(
                f"ERROR: [tool.ruff.lint].{sel_field} must be "
                 "str or list[str]",
                file=sys.stderr,
            )
            return 2
        effective_select |= {_strip_code(c) for c in codes}

    if (effective_select
            and "ALL" not in effective_select
            and "F" not in effective_select):
        print(
            "ERROR: [tool.ruff.lint] select union does not "
            "include 'F' (policy)",
            file=sys.stderr,
        )
        return 1

    # Guard 0b: all exclude variants are forbidden.
    for exc_field, exc_table in [
        ("exclude", ruff),
        ("extend-exclude", ruff),
        ("force-exclude", ruff),
        ("exclude", lint),
        ("extend-exclude", lint),
    ]:
        if exc_table.get(exc_field) is not None:
            print(
                f"ERROR: {exc_field} is forbidden by policy; "
                 "use per-file-ignores instead",
                file=sys.stderr,
            )
            return 1

    # Guard 0c: extend-per-file-ignores is forbidden.
    if lint.get("extend-per-file-ignores") is not None:
        print(
            "ERROR: extend-per-file-ignores is forbidden "
            "by policy; use per-file-ignores instead",
            file=sys.stderr,
        )
        return 1

    # Guard 1: global ignore / extend-ignore cannot contain F*.
    for field in ("ignore", "extend-ignore"):
        val = lint.get(field)
        if val is None:
            continue
        codes, ok = _iter_codes(val)
        if not ok:
            print(
                f"ERROR: [tool.ruff.lint].{field} must be "
                 "str or list[str]",
                file=sys.stderr,
            )
            return 2
        bad = [_strip_code(c) for c in codes
               if _strip_code(c).startswith(FORBIDDEN_PREFIX)]
        if bad:
            print(
                f"ERROR: global {field} contains forbidden "
                f"F* ignores: {sorted(bad)}",
                file=sys.stderr,
            )
            return 1

    # Guard 2: legacy per-file-ignores cannot contain F*.
    per_file = lint.get("per-file-ignores", {})
    if not isinstance(per_file, dict):
        print(
            "ERROR: [tool.ruff.lint.per-file-ignores] "
            "is not a table",
            file=sys.stderr,
        )
        return 2

    violations: list[str] = []
    for k, v in per_file.items():
        if not isinstance(k, str):
            continue
        if not _is_legacy_key(k):
            continue

        codes, ok = _iter_codes(v)
        if not ok and not isinstance(v, (str, list)):
            print(
                f"ERROR: per-file-ignores under {k} "
                 "must be str or list[str]",
                file=sys.stderr,
            )
            return 2
        if not ok:
            print(
                f"ERROR: per-file-ignores under {k} "
                 "contains non-string entries",
                file=sys.stderr,
            )
            return 2

        for code in codes:
            c = _strip_code(code)
            if c.startswith(FORBIDDEN_PREFIX):
                violations.append(f"{k}: {c}")

    if violations:
        print("ERROR: Forbidden F* per-file-ignores "
              "in legacy bucket:")
        for line in sorted(violations):
            print(f"  - {line}")
        return 1

    return 0


if __name__ == "__main__":
    raise SystemExit(main())
\end{lstlisting}

\section{Ruff Configuration}
\label{app:ruff}

\begin{lstlisting}[style=tomlstyle,caption={\texttt{pyproject.toml} (23 lines).},label={lst:pyproject-full}]
[tool.ruff]
target-version = "py311"
line-length = 100

[tool.ruff.lint]
select = ["E", "F", "I"]
# E402: imports after canonical docstring.
# E741: ambiguous var names in test helpers.
ignore = ["E402", "E741"]
# Keep defaults lean; do NOT globally ignore F401/F841/F541.

[tool.ruff.lint.per-file-ignores]
# CI-critical zone fill script: allow long docstrings only.
"fill_zones.py" = ["E501"]
# Generated/legacy utility scripts: cosmetic tolerances only.
"gen_*.py" = ["E501", "E401", "I001"]
"verify_*.py" = ["E501", "E401", "I001"]
# Tests: allow long lines (fixture string literals).
"test_*.py" = ["E501"]

[tool.ruff.lint.isort]
known-local-folder = ["parse_drc"]
\end{lstlisting}

\end{document}
